% Options for packages loaded elsewhere
\PassOptionsToPackage{unicode}{hyperref}
\PassOptionsToPackage{hyphens}{url}
\documentclass[
]{article}
\usepackage{xcolor}
\usepackage{amsmath,amssymb}
\setcounter{secnumdepth}{-\maxdimen} % remove section numbering
\usepackage{iftex}
\ifPDFTeX
  \usepackage[T1]{fontenc}
  \usepackage[utf8]{inputenc}
  \usepackage{textcomp} % provide euro and other symbols
\else % if luatex or xetex
  \usepackage{unicode-math} % this also loads fontspec
  \defaultfontfeatures{Scale=MatchLowercase}
  \defaultfontfeatures[\rmfamily]{Ligatures=TeX,Scale=1}
\fi
\usepackage{lmodern}
\ifPDFTeX\else
  % xetex/luatex font selection
\fi
% Use upquote if available, for straight quotes in verbatim environments
\IfFileExists{upquote.sty}{\usepackage{upquote}}{}
\IfFileExists{microtype.sty}{% use microtype if available
  \usepackage[]{microtype}
  \UseMicrotypeSet[protrusion]{basicmath} % disable protrusion for tt fonts
}{}
\makeatletter
\@ifundefined{KOMAClassName}{% if non-KOMA class
  \IfFileExists{parskip.sty}{%
    \usepackage{parskip}
  }{% else
    \setlength{\parindent}{0pt}
    \setlength{\parskip}{6pt plus 2pt minus 1pt}}
}{% if KOMA class
  \KOMAoptions{parskip=half}}
\makeatother
\usepackage{color}
\usepackage{fancyvrb}
\newcommand{\VerbBar}{|}
\newcommand{\VERB}{\Verb[commandchars=\\\{\}]}
\DefineVerbatimEnvironment{Highlighting}{Verbatim}{commandchars=\\\{\}}
% Add ',fontsize=\small' for more characters per line
\newenvironment{Shaded}{}{}
\newcommand{\AlertTok}[1]{\textcolor[rgb]{1.00,0.00,0.00}{\textbf{#1}}}
\newcommand{\AnnotationTok}[1]{\textcolor[rgb]{0.38,0.63,0.69}{\textbf{\textit{#1}}}}
\newcommand{\AttributeTok}[1]{\textcolor[rgb]{0.49,0.56,0.16}{#1}}
\newcommand{\BaseNTok}[1]{\textcolor[rgb]{0.25,0.63,0.44}{#1}}
\newcommand{\BuiltInTok}[1]{\textcolor[rgb]{0.00,0.50,0.00}{#1}}
\newcommand{\CharTok}[1]{\textcolor[rgb]{0.25,0.44,0.63}{#1}}
\newcommand{\CommentTok}[1]{\textcolor[rgb]{0.38,0.63,0.69}{\textit{#1}}}
\newcommand{\CommentVarTok}[1]{\textcolor[rgb]{0.38,0.63,0.69}{\textbf{\textit{#1}}}}
\newcommand{\ConstantTok}[1]{\textcolor[rgb]{0.53,0.00,0.00}{#1}}
\newcommand{\ControlFlowTok}[1]{\textcolor[rgb]{0.00,0.44,0.13}{\textbf{#1}}}
\newcommand{\DataTypeTok}[1]{\textcolor[rgb]{0.56,0.13,0.00}{#1}}
\newcommand{\DecValTok}[1]{\textcolor[rgb]{0.25,0.63,0.44}{#1}}
\newcommand{\DocumentationTok}[1]{\textcolor[rgb]{0.73,0.13,0.13}{\textit{#1}}}
\newcommand{\ErrorTok}[1]{\textcolor[rgb]{1.00,0.00,0.00}{\textbf{#1}}}
\newcommand{\ExtensionTok}[1]{#1}
\newcommand{\FloatTok}[1]{\textcolor[rgb]{0.25,0.63,0.44}{#1}}
\newcommand{\FunctionTok}[1]{\textcolor[rgb]{0.02,0.16,0.49}{#1}}
\newcommand{\ImportTok}[1]{\textcolor[rgb]{0.00,0.50,0.00}{\textbf{#1}}}
\newcommand{\InformationTok}[1]{\textcolor[rgb]{0.38,0.63,0.69}{\textbf{\textit{#1}}}}
\newcommand{\KeywordTok}[1]{\textcolor[rgb]{0.00,0.44,0.13}{\textbf{#1}}}
\newcommand{\NormalTok}[1]{#1}
\newcommand{\OperatorTok}[1]{\textcolor[rgb]{0.40,0.40,0.40}{#1}}
\newcommand{\OtherTok}[1]{\textcolor[rgb]{0.00,0.44,0.13}{#1}}
\newcommand{\PreprocessorTok}[1]{\textcolor[rgb]{0.74,0.48,0.00}{#1}}
\newcommand{\RegionMarkerTok}[1]{#1}
\newcommand{\SpecialCharTok}[1]{\textcolor[rgb]{0.25,0.44,0.63}{#1}}
\newcommand{\SpecialStringTok}[1]{\textcolor[rgb]{0.73,0.40,0.53}{#1}}
\newcommand{\StringTok}[1]{\textcolor[rgb]{0.25,0.44,0.63}{#1}}
\newcommand{\VariableTok}[1]{\textcolor[rgb]{0.10,0.09,0.49}{#1}}
\newcommand{\VerbatimStringTok}[1]{\textcolor[rgb]{0.25,0.44,0.63}{#1}}
\newcommand{\WarningTok}[1]{\textcolor[rgb]{0.38,0.63,0.69}{\textbf{\textit{#1}}}}
\usepackage{longtable,booktabs,array}
\newcounter{none} % for unnumbered tables
\usepackage{calc} % for calculating minipage widths
% Correct order of tables after \paragraph or \subparagraph
\usepackage{etoolbox}
\makeatletter
\patchcmd\longtable{\par}{\if@noskipsec\mbox{}\fi\par}{}{}
\makeatother
% Allow footnotes in longtable head/foot
\IfFileExists{footnotehyper.sty}{\usepackage{footnotehyper}}{\usepackage{footnote}}
\makesavenoteenv{longtable}
\setlength{\emergencystretch}{3em} % prevent overfull lines
\providecommand{\tightlist}{%
  \setlength{\itemsep}{0pt}\setlength{\parskip}{0pt}}
\usepackage{bookmark}
\IfFileExists{xurl.sty}{\usepackage{xurl}}{} % add URL line breaks if available
\urlstyle{same}
% fallback for those not using the hyperref driver hyperxmp:
\makeatletter
\@ifundefined{xmpquote}{\newcommand{\xmpquote}[1]{#1}}{}
\makeatother
\hypersetup{
  pdftitle={Trimming the Block-Scale Tax in a SpaceMiT X60 IME int8\textbackslash timesint4 GEMM Kernel: a correct{,} bit-exact optimization with no end-to-end uplift},
  pdfauthor={opensolvers/benchmarks --- IME kernel study},
  hidelinks,
  pdfcreator={LaTeX via pandoc}}

\title{Trimming the Block-Scale Tax in a SpaceMiT X60 IME
\texttt{int8}\(\times\)\texttt{int4} GEMM Kernel: a correct, bit-exact
optimization with no end-to-end uplift}
\author{opensolvers/benchmarks --- IME kernel study}
\date{July 2026}

\begin{document}
\maketitle
\begin{abstract}
llama.cpp ships a hand-written matrix-extension (IME) backend for the
SpaceMiT X60 (Ky X60) RISC-V core, built on the vendor
\texttt{smt.vmadot} instruction. Its block-scaled Q4\_0 GEMM kernel
(\texttt{gemm\_kernel\_i8i4}) runs at roughly \(0.6\times\) the
throughput of an equivalent pure-integer \texttt{s8s8s32} kernel using
the same \texttt{smt.vmadot} peak. We attribute that gap on real
silicon: stripping the per-block floating-point reduction shows the tax
is \textasciitilde31--37 \%, and it is \emph{dominated by the per-block
scale build}, not the integer\(\to\)float convert and accumulate.
Because throughput falls off almost linearly with removed instructions,
the X60 evidently does \textbf{not} overlap the IME pipe with the RVV
floating-point pipe, so cutting scale-build instruction count directly
buys time. We rewrite the scale build to use eight \texttt{vfmul.vv}
against two prebuilt row-scale vectors instead of sixteen masked
\texttt{vfmul.vf} plus four vector moves, preserving the exact
accumulator layout. The result is \textbf{bit-identical} to the stock
kernel (verified by a signature A/B on the board) and
\textbf{\textasciitilde5 \% faster} on both \texttt{M4} prefill paths.
Yet an end-to-end \texttt{llama-bench} evaluation of Qwen2.5-0.5B on the
X60 shows \textbf{no measurable prompt-processing uplift}: the prefill
throughput is statistically tied across thirteen interleaved A/B rounds,
because per-run contention noise (±15--20 \%) dwarfs the
\textasciitilde3--4 \% an isolated GEMM speedup can contribute to a
prefill that is only partly matmul. This is, deliberately, a
\emph{negative result at the application level} accompanied by a
\emph{positive, reproducible one at the kernel level} --- and an
argument for reporting both.
\end{abstract}

\section{1. Introduction}\label{introduction}

The SpaceMiT X60 (``Ky X60'', the core in the K1/M1 SoC and the Banana
Pi / Orange Pi RV2 boards) is an eight-core RISC-V application processor
implementing the RVV 1.0 vector extension at \texttt{VLEN=256}, plus a
vendor \textbf{Integer Matrix Extension (IME)} exposed through the
custom instruction \texttt{smt.vmadot}. A single \texttt{smt.vmadot}
performs a \(4\times4\) \texttt{int32} tile update from two \(4\times8\)
\texttt{int8} operands, and is the throughput engine behind the ``2 TOPS
AI'' figure quoted for the part.

llama.cpp is, at the time of writing, the only mainstream inference
stack that ships an X60 IME backend (\texttt{ggml-cpu/spacemit}, guarded
by \texttt{GGML\_CPU\_RISCV64\_SPACEMIT}). Its central kernel,
\texttt{spacemit\_kernels::ime1::gemm\_kernel\_i8i4}, is a
\emph{block-scaled} Q4\_0 matrix multiply: \texttt{int8} activations
times 4-bit weights, with a per-32-element \texttt{fp16} scale,
accumulated in \texttt{fp32}. A companion microbenchmark in this
repository (\texttt{ime/}) established that this block-scaled kernel
tops out around \(0.6\times\) the throughput of a pure \texttt{s8s8s32}
kernel built on the same \texttt{smt.vmadot} peak. The difference is the
\emph{block-scale tax}: work the integer kernel never does --- one
\texttt{int32}\(\to\)\texttt{fp32} convert and one scale-multiply per
32-element K-block, per output tile.

Prior notes speculated that the tax ``needs a kernel rewrite, not a
driver tweak.'' This paper does the rewrite, measures it rigorously on
the board, and reports the outcome honestly. The kernel-level result is
a clean, bit-exact \textasciitilde5 \%; the application-level result is
\emph{nothing measurable}. We consider the second finding as important
as the first: it is a concrete, root-caused example of a microbenchmark
win that does not survive contact with a real workload on real, shared
hardware, and it explains precisely why.

\section{2. Background}\label{background}

\subsection{2.1 The IME kernel path}\label{the-ime-kernel-path}

\texttt{gemm\_kernel\_i8i4} dispatches on the number of activation rows:

\begin{itemize}
\tightlist
\item
  \textbf{\texttt{M4}} (\texttt{count\_m} \(\geq 4\)) --- the
  register-blocked path used by \emph{prefill} (\texttt{pp}, prompt
  processing). It computes a \(4\times16\) \texttt{fp32} output tile per
  call and dominates prompt-processing compute. It has two variants,
  with and without a per-block zero point (\texttt{M4} no-ZP is the one
  Q4\_0 uses; \texttt{M4} ZP serves zero-point quant formats).
\item
  \textbf{\texttt{M1}} (\texttt{count\_m\ \textless{}\ 4}) --- the GEMV
  path used by \emph{decode} (\texttt{tg}, token generation), one
  activation row against the weight matrix.
\end{itemize}

Per K-block, the \texttt{M4} path issues sixteen \texttt{smt.vmadot}
(the matmul), then a floating-point \emph{reduction} consisting of (i) a
\textbf{scale build} (\texttt{LOAD\_SCALE\_4x16\_FP16}) that loads
sixteen \texttt{fp16} weight scales, widens them, and folds in the four
per-row activation scales, and (ii) a \textbf{convert + accumulate} (one
\texttt{vfcvt.f.x.v} and one \texttt{vfmacc.vv}, both at
\texttt{LMUL=8}) that turns the \texttt{int32} tile accumulator into
scaled \texttt{fp32} and adds it into the running output. The integer
accumulator must reset every K-block because each Q4\_0 block carries
its own scale --- this is the irreducible structural reason the
reduction cannot be hoisted across blocks.

\subsection{2.2 The accumulator layout (why the scale build is
awkward)}\label{the-accumulator-layout-why-the-scale-build-is-awkward}

\texttt{smt.vmadot} writes a \(4\times4\) tile into an even/odd register
pair. The four accumulators \texttt{v16/v18/v20/v22} therefore hold four
\emph{column-groups}, each \(4\ \text{rows}\times4\ \text{cols}\), and
the store macro \texttt{SAVE\_RESULT\_4x16} reveals the exact lane
mapping: within each \texttt{fp32} output register, the low four lanes
are one output row and the high four lanes are the next row, for a fixed
column-group. The combined scale for that register is thus
\(A_{\text{scale}}[\text{row}]\cdot B_{\text{scale}}[\text{col}]\) with
the \emph{weight} scale replicated across the two row-halves and the two
\emph{activation} scales applied per-half. The stock code realises this
with overlapping masked loads (to replicate the weight scales) followed
by sixteen predicated \texttt{vfmul.vf} (eight at \texttt{vl=4}, eight
at \texttt{vl=8} under a \texttt{0xf0} mask) to apply the row scales ---
plus four \texttt{vmv.v.v} copies and several \texttt{vsetvli}
(vector-type) switches.

\section{3. The block-scale tax,
measured}\label{the-block-scale-tax-measured}

All measurements are single-core on an Orange Pi RV2 (X60 @ 1.6 GHz,
\texttt{performance} governor), pinned to core 0 of cluster 0, reported
as the \emph{clean peak} (the maximum over \(\geq 4\) process launches;
contention only ever lowers a run, so the maximum is the true
single-core capability). Throughput is single-core GOP/s (\(2MNK/t\)).

To localize the tax we build variants of the \texttt{M4} no-ZP path that
strip parts of the per-block reduction (timing-only; these variants are
numerically wrong by construction):

\begin{itemize}
\tightlist
\item
  \textbf{\texttt{-scale-build}}: remove
  \texttt{LOAD\_SCALE\_4x16\_FP16}, keep the convert + accumulate.
\item
  \textbf{\texttt{-all\ FP} (vmadot only)}: remove the entire reduction,
  leaving only the sixteen \texttt{smt.vmadot} and the loop overhead.
\end{itemize}

{\def\LTcaptype{none} % do not increment counter
\begin{longtable}[]{@{}
  >{\raggedright\arraybackslash}p{(\linewidth - 6\tabcolsep) * \real{0.2500}}
  >{\raggedleft\arraybackslash}p{(\linewidth - 6\tabcolsep) * \real{0.2500}}
  >{\raggedleft\arraybackslash}p{(\linewidth - 6\tabcolsep) * \real{0.2500}}
  >{\raggedleft\arraybackslash}p{(\linewidth - 6\tabcolsep) * \real{0.2500}}@{}}
\toprule\noalign{}
\begin{minipage}[b]{\linewidth}\raggedright
\(M\times N\times K\)
\end{minipage} & \begin{minipage}[b]{\linewidth}\raggedleft
stock
\end{minipage} & \begin{minipage}[b]{\linewidth}\raggedleft
\(-\)scale-build
\end{minipage} & \begin{minipage}[b]{\linewidth}\raggedleft
\(-\)all FP (vmadot only)
\end{minipage} \\
\midrule\noalign{}
\endhead
\bottomrule\noalign{}
\endlastfoot
\(512^3\) & 22.5 & 29.6 & 32.7 \\
\(768^3\) & 27.5 & 36.7 & 41.3 \\
\(512\times512\times2048\) & 27.5 & 38.7 & 43.6 \\
\end{longtable}
}

Two things follow. First, the full floating-point tax is large ---
31--37 \% --- and the \texttt{vmadot}-only ceiling at \(768^3\) (41.3)
and large K (43.6) approaches the pure \texttt{s8s8s32} ceiling
(\textasciitilde42), confirming the matmul engine itself is not the
bottleneck. Second, and decisively: removing \textasciitilde40
scale-build instructions recovers \textasciitilde31 \%, whereas removing
the \textasciitilde16-instruction convert + accumulate recovers only
\textasciitilde10 \%. The recovered time is roughly proportional to
removed instruction count. \textbf{The X60 does not appreciably overlap
the IME pipe with RVV floating-point} --- the two phases are effectively
serial --- so the tax is paid in issue slots, and cutting instruction
count in the dominant phase (the scale build) is the lever.

\section{4. Optimization: a register-efficient scale
build}\label{optimization-a-register-efficient-scale-build}

The scale build produces eight \texttt{fp32} registers \texttt{v8..v15},
each holding \(B_{\text{scale}}[\text{col-group}]\) replicated across
both row-halves and scaled by two activation scales. The stock code
applies the row scales with sixteen predicated \texttt{vfmul.vf} (half
of them at \texttt{vl=8} under a mask, i.e.~doing four lanes of work in
eight lanes of time) plus four \texttt{vmv.v.v} and repeated
\texttt{vsetvli} switches.

We instead precompute two \emph{row-scale vectors} once per K-block ---
\(\mathbf{a}_{\text{even}}=[A_0^{\times4}, A_1^{\times4}]\) and
\(\mathbf{a}_{\text{odd}}=[A_2^{\times4}, A_3^{\times4}]\) --- with two
\texttt{vfmv.v.f} and two \texttt{vfmerge.vfm}, then form the eight
combined-scale registers with eight full \texttt{vfmul.vv}:

\begin{Shaded}
\begin{Highlighting}[]
\NormalTok{    vfmv}\OperatorTok{.}\NormalTok{v}\OperatorTok{.}\NormalTok{f     v2}\OperatorTok{,}\NormalTok{ f1            }\CommentTok{; a\_even = [A0,A0,A0,A0, A1,A1,A1,A1]}
\NormalTok{    vfmerge}\OperatorTok{.}\NormalTok{vfm  v2}\OperatorTok{,}\NormalTok{ v2}\OperatorTok{,}\NormalTok{ f2}\OperatorTok{,}\NormalTok{ v0    }\CommentTok{;   (v0 = 0xf0 mask {-}\textgreater{} high half)}
\NormalTok{    vfmv}\OperatorTok{.}\NormalTok{v}\OperatorTok{.}\NormalTok{f     v3}\OperatorTok{,}\NormalTok{ f3            }\CommentTok{; a\_odd  = [A2,A2,A2,A2, A3,A3,A3,A3]}
\NormalTok{    vfmerge}\OperatorTok{.}\NormalTok{vfm  v3}\OperatorTok{,}\NormalTok{ v3}\OperatorTok{,}\NormalTok{ f4}\OperatorTok{,}\NormalTok{ v0}
\NormalTok{    vfmul}\OperatorTok{.}\NormalTok{vv     v8}\OperatorTok{,}\NormalTok{  v4}\OperatorTok{,}\NormalTok{ v2       }\CommentTok{; B(cg0) * a\_even}
\NormalTok{    vfmul}\OperatorTok{.}\NormalTok{vv     v9}\OperatorTok{,}\NormalTok{  v4}\OperatorTok{,}\NormalTok{ v3       }\CommentTok{; B(cg0) * a\_odd}
\NormalTok{    vfmul}\OperatorTok{.}\NormalTok{vv     v10}\OperatorTok{,}\NormalTok{ v5}\OperatorTok{,}\NormalTok{ v2       }\CommentTok{; ... four column{-}groups x \{even,odd\}}
\NormalTok{    ...}
\NormalTok{    vfmul}\OperatorTok{.}\NormalTok{vv     v15}\OperatorTok{,}\NormalTok{ v7}\OperatorTok{,}\NormalTok{ v3}
\end{Highlighting}
\end{Shaded}

This drops the four \texttt{vmv.v.v}, halves the multiply count (16
masked \texttt{vfmul.vf} \(\to\) 8 full \texttt{vfmul.vv}), and removes
the \texttt{vl=4}/\texttt{vl=8} \texttt{vsetvli} alternation. The
\texttt{v8..v15} layout is byte-for-byte the layout the stock kernel
produced, so the downstream \texttt{vfcvt}/\texttt{vfmacc} and
\texttt{SAVE\_RESULT\_4x16} are unchanged and the numeric result is
identical. The change is \textasciitilde40 lines confined to a single
macro (patch: \texttt{ime/llama-ime1-scalebuild-opt.patch}).

Both \texttt{M4} variants (no-ZP and ZP) share
\texttt{LOAD\_SCALE\_4x16\_FP16}, so the same macro applies to both; the
swap is register-safe in the ZP path (it preserves \texttt{v1}, the
zero-point gather-index register, and clobbers only registers the stock
macro already clobbers).

\section{5. Methodology}\label{methodology}

\textbf{Build.} The kernel is edited on an x86-64 host and
cross-compiled with the SpaceMiT GCC toolchain into
\texttt{libggml-cpu.so}; an incremental rebuild is \textasciitilde20 s.
The shared object is copied to the board, where the microbenchmark and
\texttt{llama-bench} load it via \texttt{rpath}.

\textbf{Correctness.} The microbenchmark harness
(\texttt{ime/bench\_i8i4.cpp}) gained a \texttt{chk} mode that fills
both operand buffers with varied bytes constrained to
\texttt{{[}1,0x3f{]}} --- guaranteeing every \texttt{fp16}/\texttt{fp32}
value the kernel might read is finite and normal, while still varying
per (row, column, block) so a wrong scale or tile layout surfaces as a
mismatch. It then prints a signature (sum, sum-of-squares, max of the
output). We A/B the \emph{modified} kernel's signature against the
\emph{stock} kernel's on byte-identical inputs; a rewrite is accepted
only if the signatures match exactly. Throughput is data-independent
(straight-line integer and floating-point, no data branches), so timing
is measured separately with representative data. The \texttt{chk} mode
also gained \texttt{zp} and \texttt{m1} selectors so all three kernel
paths can be exercised.

\section{6. Kernel results}\label{kernel-results}

Signatures matched exactly for every size on both \texttt{M4} paths ---
the optimization is \textbf{bit-identical} to the stock kernel.
Clean-peak throughput (paired A/B, single core):

{\def\LTcaptype{none} % do not increment counter
\begin{longtable}[]{@{}
  >{\raggedright\arraybackslash}p{(\linewidth - 6\tabcolsep) * \real{0.2500}}
  >{\raggedright\arraybackslash}p{(\linewidth - 6\tabcolsep) * \real{0.2500}}
  >{\raggedleft\arraybackslash}p{(\linewidth - 6\tabcolsep) * \real{0.2500}}
  >{\raggedleft\arraybackslash}p{(\linewidth - 6\tabcolsep) * \real{0.2500}}@{}}
\toprule\noalign{}
\begin{minipage}[b]{\linewidth}\raggedright
kernel path
\end{minipage} & \begin{minipage}[b]{\linewidth}\raggedright
used by
\end{minipage} & \begin{minipage}[b]{\linewidth}\raggedleft
\(512^3\)
\end{minipage} & \begin{minipage}[b]{\linewidth}\raggedleft
\(768^3\)
\end{minipage} \\
\midrule\noalign{}
\endhead
\bottomrule\noalign{}
\endlastfoot
\texttt{M4} no-zero-point & Q4\_0 prefill & \(23.0\to24.1\) (+4.8 \%) &
\(27.6\to29.2\) (+5.8 \%) \\
\texttt{M4} zero-point & ZP-quant prefill & \(20.6\to21.7\) (+5.3 \%) &
\(24.7\to26.0\) (+5.3 \%) \\
\texttt{M1} / GEMV & decode (\texttt{tg}) & 12.1 --- N/A & --- \\
\end{longtable}
}

The \texttt{M4} gain is a consistent \textasciitilde5 \% across
\(512^3\)--\(2048^2\times512\). The \texttt{M1} / GEMV path is
\textbf{not applicable}: its scale build is already minimal (four
\texttt{vfmul.vf} for a single activation row, no masking, no
\texttt{vmv}), so there is no sixteen-\texttt{vfmul} tax to cut, and
GEMV on this part is memory-bound anyway (12.1 GOP/s, roughly half the
\texttt{M4} rate --- the weight stream, used once per token, dominates).

\section{7. End-to-end evaluation}\label{end-to-end-evaluation}

We ran \texttt{llama-bench} on Qwen2.5-0.5B (Q4\_0) with the stock and
patched \texttt{libggml-cpu.so}, four threads pinned to cluster 0
(\texttt{taskset\ -c\ 0-3}), comparing the two shared objects in the
same session. \texttt{pp512} (prefill) exercises the optimized
\texttt{M4} GEMM; \texttt{tg128} (decode) exercises the untouched
\texttt{M1} path and serves as a negative control.

\textbf{Correctness holds end-to-end.} The patched build produces
coherent, correct output (\emph{``The capital of France is Paris.''})
and \texttt{tg128} is bit-stable at \textbf{7.25 t/s} on both builds ---
as expected, since decode never touches the changed code.

\textbf{Prefill shows no measurable uplift.} Over thirteen interleaved
A/B rounds:

{\def\LTcaptype{none} % do not increment counter
\begin{longtable}[]{@{}lrr@{}}
\toprule\noalign{}
metric & stock & patched \\
\midrule\noalign{}
\endhead
\bottomrule\noalign{}
\endlastfoot
\texttt{pp512} mean (t/s) & 79.8 & 79.5 \\
\texttt{pp512} median (t/s) & 82.3 & 79.9 \\
\texttt{pp512} peak (t/s) & \textbf{91.3} & 86.3 \\
\texttt{pp512} range (t/s) & 62.8--91.3 & 67.0--86.3 \\
\texttt{tg128} (t/s) & 7.25 ± 0.01 & 7.25 ± 0.01 \\
\end{longtable}
}

The two prompt-processing distributions are statistically
indistinguishable --- means within 0.3 \%, and if anything the stock
peak and median are marginally higher (noise). The cause is visible in
the range: \texttt{pp512} swings ±15--20 \% run-to-run, a
\textasciitilde45 \% total spread, because prefill is memory- and
L2-bandwidth-heavy and this is a shared, multi-tenant part --- the four
cluster-0 cores share one 512 KB L2, and buffer placement additionally
triggers the \texttt{malloc}-dependent cache-set aliasing documented in
the companion \texttt{ime/} study. That noise floor is an order of
magnitude larger than the uplift a 5 \% \emph{kernel} speedup can
contribute to a prefill that is only partly matmul (attention, RoPE,
softmax, norms, and the sampling path also cost). The expected
end-to-end gain, \textasciitilde3--4 \%, is simply not resolvable here.

\section{8. Discussion}\label{discussion}

\textbf{Why the tax is otherwise fundamental.} The convert and
accumulate (\texttt{vfcvt}/\texttt{vfmacc}) and the weight-scale loads
are near-irreducible, and the per-block scale cannot be deferred across
K-blocks because each Q4\_0 block has its own scale. The one lever that
could hide the \emph{whole} floating-point phase --- software-pipelining
the next block's \texttt{smt.vmadot} under the current block's
reduction, given that the two pipes are serial --- is
\textbf{register-blocked}. A whole \(4\times16\) tile needs eight
accumulator registers; double-buffering it needs sixteen, but the
\texttt{fp32} output (\texttt{v24..v31}), the \texttt{int32} accumulator
(\texttt{v16..v23}), and the scale vectors (\texttt{v8..v15}) already
fill the 32-register file.

\textbf{Why offloading the reduction to the idle cores does not help.}
The four non-IME cores (cluster 1) sit idle during IME GEMM, which
invites moving the floating-point reduction there. It does not pay: the
two clusters have \emph{separate} 512 KB L2 caches, and block scaling
forbids accumulating \texttt{int32} across blocks, so an IME core would
have to ship a per-block \texttt{int32} partial tile to a floating-point
core --- \(M\cdot N\cdot K/8\) bytes in total (\textasciitilde16 MB at
\(512^3\) versus a 1 MB output), across the inter-cluster interconnect,
to save arithmetic that is only \textasciitilde0.3 flop/byte. The
intermediate lives in vector registers today at zero memory traffic;
externalizing it is a net loss. The spare cores are worth more running
an independent RVV GEMM --- and indeed a well-threaded RVV build already
beats IME end-to-end on this part, because IME is confined to cluster 0.

\textbf{Portability.} The optimization is specific to the SpaceMiT IME
microkernel, but the packing it feeds (a pre-transposed
\(B\)-is-\(N\times K\) tile layout, the same shape \texttt{mmt4d} and
ggml-repack produce) is the layout any framework matrix-unit backend
expects. The X60 IME microkernel is a building block that MLAS (hence
ONNX Runtime), XNNPACK (TensorFlow Lite, PyTorch mobile), and ruy
(TensorFlow Lite) do not ship for this hardware; the block-scale
accounting here transfers directly, as does the conclusion that
scale-build instruction count --- not the matrix instruction --- is
where a serial-pipe matrix core spends the quantization overhead.

\textbf{Where the win \emph{would} land.} A quieter or single-tenant
board would surface the \textasciitilde3--4 \% prefill gain; more
usefully, a part whose matrix unit spans more than one cluster's cores
(lifting the cluster-0 cap) would let IME win end-to-end, at which point
a 5 \% kernel improvement is 5 \% of a path that matters.

\section{9. Conclusion}\label{conclusion}

We closed the largest reducible part of the SpaceMiT X60 IME block-scale
tax with a register-efficient scale build: eight \texttt{vfmul.vv} and
two prebuilt row-scale vectors in place of sixteen masked
\texttt{vfmul.vf}, four vector moves, and their \texttt{vsetvli} churn.
The change is bit-identical to the stock kernel and \textasciitilde5 \%
faster on both \texttt{M4} prefill paths, verified on real X60 silicon.
It buys \textbf{no measurable end-to-end speedup} on this board, because
GEMM is only part of prefill and the application-level measurement noise
(±15--20 \%) is far larger than the achievable gain. Both results are
real, and reporting only the first would be misleading. The kernel
improvement is free, correct, and upstreamable; its practical value is
gated by the microarchitecture (serial IME/FP pipes, register pressure)
and the platform (a four-core, shared-L2 IME cluster), not by the
kernel.

\section{Artifact availability}\label{artifact-availability}

The microbenchmark harness (\texttt{ime/bench\_i8i4.cpp}, with the
\texttt{chk}/\texttt{zp}/\texttt{m1} A/B modes), the
pure-\texttt{s8s8s32} reference (\texttt{ime/}), and the optimization
itself (\texttt{ime/llama-ime1-scalebuild-opt.patch}) are in this
repository. All numbers were produced on an Orange Pi RV2 (SpaceMiT
K1/X60) at 1.6 GHz, \texttt{performance} governor, single core 0 for the
microbenchmarks and cluster 0 (\texttt{-t4}) for \texttt{llama-bench}.

\section{References}\label{references}

\begin{enumerate}
\def\labelenumi{\arabic{enumi}.}
\tightlist
\item
  llama.cpp / ggml --- SpaceMiT IME backend (\texttt{ggml-cpu/spacemit},
  \texttt{GGML\_CPU\_RISCV64\_SPACEMIT}; \texttt{ime1\_kernels.cpp}).
\item
  SpaceMiT K1/M1 (Ky X60) technical reference --- IME
  \texttt{smt.vmadot}, \texttt{xsmtvdot} assembler support (binutils
  \(\geq\) 2.46).
\item
  RISC-V ``V'' Vector Extension, version 1.0.
\item
  Microsoft MLAS \texttt{MlasSQNBitGemm} (the \texttt{SQ4Bit…CompInt8}
  quantized-GEMM scheme the llama.cpp kernel names derive from).
\item
  \texttt{opensolvers/benchmarks} --- \texttt{ime/README.md} (pure
  \texttt{s8s8s32} ceiling, cache-set aliasing) and \texttt{ime-llama}
  study notes.
\end{enumerate}

\end{document}
